% Generated draft: requires mathematical and source review.
\documentclass[11pt]{article}
\usepackage[margin=25mm]{geometry}
\usepackage{fontspec}
\setmainfont{Latin Modern Roman}
\usepackage{amsmath,amssymb,mathtools,unicode-math}
\setmathfont{Latin Modern Math}
\usepackage{enumitem,needspace,xurl,hyperref}
\setlength{\parindent}{0pt}
\setlength{\parskip}{5pt}
\setlength{\emergencystretch}{3em}
\newcommand{\N}{\mathbb{N}}
\newcommand{\Z}{\mathbb{Z}}
\newcommand{\Q}{\mathbb{Q}}
\newcommand{\R}{\mathbb{R}}
\newcommand{\C}{\mathbb{C}}
\newcommand{\F}{\mathbb{F}}
\newcommand{\A}{\mathbb{A}}
\newcommand{\PP}{\mathbb{P}}
\newcommand{\E}{\mathbb{E}}
\newcommand{\eps}{\varepsilon}
\newcommand{\dd}{\,\mathrm{d}}
\begin{document}
\section*{UnsolvedMath ID: 30005812}
\textbf{Source problem code:} catalogue entry
\subsection{Definitions and mathematical statement}
The supplied reconstruction asks a finiteness question about regular graphs. Its graph conventions are not specified. For a provisional precise interpretation, let \(G=(V,E)\) be a finite, connected, simple, undirected graph containing a cycle. Call \(G\) regular if every vertex has the same degree \(k\); the degree \(k\) is allowed to vary between graphs. Define
\[
D(G)=\max_{u,v\in V}d_G(u,v),\qquad
g(G)=\min\{\lvert E(C)\rvert:C\text{ is a cycle in }G\},
\]
where \(d_G(u,v)\) is the length of a shortest path from \(u\) to \(v\). Write \(K_n\) for the complete graph on \(n\) vertices and \(C_m\) for the cycle on \(m\) vertices.

Under these provisional conventions, the reconstructed question is whether, for every fixed real number \(c>0\), the set
\[
\mathcal G_c=
\left\{
[G]_{\cong}\;\middle|\;
\begin{gathered}
G\text{ satisfies the conventions above and is regular},\\
g(G)=(2-c)D(G),\\
G\not\cong K_n\quad\text{for every integer }n\geq3,\\
G\not\cong C_{2n+1}\quad\text{for every integer }n\geq1
\end{gathered}
\right\}
\]
is finite. Here \([G]_{\cong}\) denotes the isomorphism class of \(G\); any finiteness bound may depend on \(c\).

This is only a conditional formalization of the supplied reconstruction. The original fragment instead displays
\[
g(G)=2D(G)-1,
\]
and asks whether their number is finite, while referring to regular Moore graphs and to complete graphs and odd cycles as exceptional families. The displayed additive relation agrees with the reconstructed relation only when \(cD(G)=1\); it therefore does not establish the same fixed-\(c\) question when diameters vary. Moreover, under the provisional conventions, complete graphs and odd cycles have \(g(G)/D(G)>2\), so their exclusion is redundant in the reconstructed relation for \(c>0\). This leaves their intended role unresolved. The fragment does not define the intended meaning of ``regular Moore graphs.'' The intended relation, graph domain, and scope of finiteness require source review before this draft can be treated as a definitive statement.
\subsection{Short English statement}
For each fixed real \(c>0\), the supplied reconstruction asks whether only finitely many regular graphs, up to isomorphism and excluding complete graphs and odd cycles, satisfy \(g(G)=(2-c)D(G)\). The original fragment instead gives \(g(G)=2D(G)-1\). The intended equation and graph conventions remain unresolved; the displayed formalization is conditional and requires human mathematical review.
\subsection{Sources}
\begin{itemize}
\item Local UnsolvedMath registry record, permanent ID \(30005812\), title ``Finiteness of Generalized Moore Graphs,'' source problem code \texttt{catalogue entry}, in the Oberwolfach Reports Open Problems collection. The supplied record associates the problem with 2024 and provides no named proposer.
\item Supplied primary report citation: \emph{Discrete Geometry} (2024), Oberwolfach Reports, pp.\ 137--202, DOI \texttt{10.4171/owr/2024/3}. \url{https://doi.org/10.4171/owr/2024/3}. No problem-specific page or named contributor is supplied.
\item Supplied literature source: ``Existence of a Moore graph of degree 57 is still open.'' \url{https://arxiv.org/abs/2210.09577}.
\item Supplied literature source: ``Moore graphs and beyond.'' \url{https://www.combinatorics.org/files/Surveys/ds14/ds14v2-2013.pdf}.
\item Literature assessment dated exactly \texttt{2026-08-22}: the supplied curation labels the problem open, states that known Moore-graph classification leaves the exceptional degree-57 diameter-two case unresolved, and states that the broader near-Moore formulation in the fragment is not settled by those classification theorems. This is supplied dated curation, not an independently verified assessment; its time-sensitive status and cited claims require independent checking.
\item The registry labels the reconstructed statement \texttt{corrected\_verified} and reports that it was checked against the official Oberwolfach report. This is a supplied verification claim, not verification performed for this draft, and it does not resolve the discrepancy between the two supplied equations.
\end{itemize}
\end{document}
